% Reconstructed 2026-11-09 from text extraction of the RECOVERED 60-page reader PDF.
% The original 37-block TeX is lost. Display mathematics, superscripts, and subscript formatting
% may have been flattened by PDF text extraction. Read DER_VI-a_reader.pdf as authoritative.
% These are not byte-identical to the original source and were not certified as proof edits.
\begin{dcomment}{context}{p1}{The axioms are being found}
Grothendieck wants to construct the homotopy theory of diagrams from categories themselves. The expression “catégorie de modèles” here initially means a category supplied with distinguished weak equivalences; no classes of Quillen cofibrations and fibrations have yet been specified. We write W\ensuremath{\infty} for functors inducing weak homotopy equivalences of nerves, and Q : Cat \ensuremath{\to}Cat[W \ensuremath{-}1] for ordinary localization. Categories and localizations are understood in fixed, sufficiently large universes. This avoids mistaking the construction of Q for a smallness theorem about its hom sets.

The changing list of axioms is part of the argument. In particular, the L4 of manuscript page 2 is the absolute assertion that an aspherical functor is a weak equivalence. The relative assertion for a map X \ensuremath{\to}Y over an arbitrary third category S appears only as L4 bis on page 44. Later proofs require this stronger formulation. The modern definition uses the relative statement from the outset; compare Cisinski (2004, Definition 1.1.3 and Remark 1.1.4). Similarly, the L5 on page 2 concerns descent of fibrewise equivalences, whereas the L5 on page 43 concerns recovery of fibrewise equivalences from a total equivalence. They are different assertions, with different strength. The interval [1] explains why these axioms can remember homotopy. A natural transformation F \ensuremath{\Rightarrow}G is a functor A \ensuremath{\times} [1] \ensuremath{\to}B with prescribed endpoint restrictions. L2 makes both endpoints inverse to the projection after localization. Thus a natural transformation becomes an equality between the induced maps of localized categories. This observation will repeatedly turn an elementary adjunction into a weak equivalence without requiring a model structure.

The optional L3 quater also has a precise later history. Under the modern relative axioms, the inclusion W\ensuremath{\infty}\ensuremath{\subseteq}W is a theorem, established by Cisinski (2004, Theorem 2.2.11). That theorem does not license replacing the initial absolute L4 by its relative form without explanation. In these pages the additional hypothesis is explicitly adopted at manuscript page 7, and we retain the distinction when tracing the proofs.

\end{dcomment}

\begin{dcomment}{exposition}{p2}{What the first fibration axiom already gives}
Grothendieck’s short implication from L1, L2 bis and the original L5 can be expanded into a useful collection of elementary arguments. It also clarifies why that first L5 becomes dispensable after relative descent has been imposed. First apply L5 to the cartesian functor [1] \ensuremath{\times} A \ensuremath{\to}A between fibrations over A. Its fibre map is [1] \ensuremath{\to}\{\ensuremath{*}\}, so L2 follows. For the projection p : A \ensuremath{\times} [1] \ensuremath{\to}A and its sections e0, e1, two-out-of-three gives e0, e1 \ensuremath{\in}W. If H : A \ensuremath{\times} [1] \ensuremath{\to}B is a natural transformation between F and G, it follows that F \ensuremath{\in}W if and only if G \ensuremath{\in}W: both are equivalent to H \ensuremath{\in}W. Moreover, Qe0 = Qp\ensuremath{-}1 = Qe1, whence QF = QG. An adjoint pair F \ensuremath{\dashv}G has GF and FG connected to the corresponding identities by its unit and counit. Both composites belong to W, and L1(c) gives F, G \ensuremath{\in}W. Applying this to the terminal-object adjunction proves L3, and the initial-object case follows as well.

There is a particularly concrete way to obtain invariance under opposites without assuming it. Let T(A) have objects the arrows a \ensuremath{\to}b of A, with a morphism from a \ensuremath{\to}b to a\ensuremath{\prime} \ensuremath{\to}b\ensuremath{\prime} given by a\ensuremath{\prime} \ensuremath{\to}a and b \ensuremath{\to}b\ensuremath{\prime} making the evident rectangle commute. The two projections

T(A)op \ensuremath{-}\ensuremath{\to}A,

T(A)op \ensuremath{-}\ensuremath{\to}Aop are fibrations. Their fibres are respectively (a\textbackslash{}A)op and A/b, each with a terminal object. L5, applied to each projection as a cartesian functor to the identity fibration, makes both projections elements of W. Their naturality in A and two-out-of-three now imply

T(f)op \ensuremath{\in}W

f op \ensuremath{\in}W. f \ensuremath{\in}W \ensuremath{\Leftarrow}\ensuremath{\Rightarrow} \ensuremath{\Leftarrow}\ensuremath{\Rightarrow} Here the fibration assertion can be checked directly: a change of the first endpoint is made by precomposition before taking opposites, and a change of the second endpoint by postcomposition. These are the specified cartesian lifts after reversal.

Now let f : X \ensuremath{\to}Y be a map over S such that every f/s : X/s \ensuremath{\to}Y/s belongs to W. The construction \ensuremath{\Psi} on manuscript pages 17--20 produces split opfibrations with these maps as their fibre maps, and \ensuremath{\Psi}(f) preserves the chosen cocartesian lifts. After taking opposites, L5 applies, so \ensuremath{\Psi}(f) \ensuremath{\in}W. The natural maps X \ensuremath{\to}\ensuremath{\Psi}(X) and Y \ensuremath{\to}\ensuremath{\Psi}(Y ) are homotopy equivalences by the adjunction exhibited below; hence two-out-of-three gives f \ensuremath{\in}W. This proves even L4 bis, and therefore L4. Each ingredient is independent of the later, much stronger L5. The argument verifies the implication recorded here while anticipating only constructions that Grothendieck gives explicitly in the next few pages.

\end{dcomment}

\begin{dcomment}{exposition}{p3}{Products and sums survive the localization}
Grothendieck asks for concrete conditions ensuring Lb and Lc because a homotopy theory should still allow objects to be assembled and compared. A useful answer is available as soon as relative L4 bis is admitted. We give the argument, including the passage from closure of weak equivalences to the universal property of products and coproducts.

Suppose W satisfies L1, L2 and L4 bis. If f : A \ensuremath{\to}B lies in W, then f \ensuremath{\times} idC lies in W for every C. View this map over C. Its comma map at c is f \ensuremath{\times} idC/c. Since C/c has a terminal object, the two projections A \ensuremath{\times} (C/c) \ensuremath{\to}A and B \ensuremath{\times} (C/c) \ensuremath{\to}B are homotopy equivalences: the terminal object defines sections and a contraction natural transformation. They are in W by the preceding block. Two-out-of-three makes the comma map a weak equivalence, and L4 bis gives the assertion. The same axiom, applied over a discrete indexing category I, shows that ` i fi \ensuremath{\in}W whenever every fi \ensuremath{\in}W. The latter closure is also Cisinski (2004, Proposition 1.1.5).

For completeness, localization of a product of relative categories is the product of their localizations. One can see this directly from the generators: (w, id) and (id, w\ensuremath{\prime}) become invertible, their commutation squares still commute after inversion, and a pair of finite zigzags can be interleaved to give a zigzag in the product. Consequently both the diagonal and the binary product functor descend. The adjunction \ensuremath{\Delta}\ensuremath{\dashv}\ensuremath{\times} descends with its original unit and counit, since those are natural transformations between relative functors and the triangle identities remain identities. Its descended right adjoint is therefore a binary product in Cat[W \ensuremath{-}1]. The same argument with ` \ensuremath{\dashv}\ensuremath{\Delta}gives binary coproducts. The adjunctions defining the initial and terminal objects descend without any closure hypothesis beyond identities belonging to W.

Thus finite products and coproducts are preserved by Q under these concrete assumptions. This is an ordinary localization argument, not a claim that arbitrary homotopy limits or colimits agree with ordinary ones. In particular, it does not supply the derived Kan extensions needed for the derivator. Those require the diagram categories whose construction begins here.

\end{dcomment}

\begin{dcomment}{exposition}{p4}{A diagram written as its category of elements}
Grothendieck passes to fibrations in order to replace an externally indexed diagram by one category carrying a projection to its index. This is the same organizing construction used in descent; the relevant background is Grothendieck (1971, Exposé VI, \S{}\S{}8--9). It is useful to write out the variance because both directions of the construction will soon be needed.

For F : Sop \ensuremath{\to}Cat, its Grothendieck construction has objects (s, x) with x \ensuremath{\in}F(s). An arrow (s, x) \ensuremath{\to}(t, y) consists of u : s \ensuremath{\to}t and a : x \ensuremath{\to}F(u)y in F(s). Composition is

(v, b) \ensuremath{\circ}(u, a) = (vu, F(u)(b) \ensuremath{\circ}a).

The chosen cartesian lift of u at (t, y) is (u, id) : (s, F(u)y) \ensuremath{\to}(t, y). The strict functoriality of F makes these choices compose strictly. Conversely, a normalized splitting of a fibration gives strict pullback functors u\ensuremath{*}, with (vu)\ensuremath{*}= u\ensuremath{*}v\ensuremath{*}and identity pullbacks equal to identities. These equations, rather than the existence of unspecified isomorphisms, are what is meant by a splitting.

A natural transformation F \ensuremath{\to}G induces a functor of the total categories that preserves the specified lifts. Conversely, such a functor restricts to the components of a strict natural transformation. Therefore the weak equivalences in the diagram category correspond exactly to fibrewise weak equivalences among splitting-preserving functors. With fibrations encoded by their fibre families, this is the identification Grothendieck calls an isomorphism. If arbitrary underlying object sets for total categories are allowed, the usual construction and reconstruction are canonically isomorphic over S, and the comparison is expressed as an equivalence. No homotopical argument depends on this bookkeeping convention.

The subsequent forgetful functor to Fib S changes the morphisms: there one allows all functors over S. Forgetting a splitting is therefore not merely forgetting redundant labels. The natural comparison maps used later must be checked in the category whose morphisms they actually preserve. This is the source of the obstruction on manuscript pages 32--39, and of the strictification work beginning on page 52.

\end{dcomment}

\begin{dcomment}{examples}{p7}{A local system in the homotopy category misses twisting}
Grothendieck’s passage from F to F\textasciimacron{} : \ensuremath{\Pi}(Sop) \ensuremath{\to}HotW records monodromy, but he also signals on page 5 that richer coherence remains. A familiar pair of bundles makes the difference visible. We use ordinary spaces to display the phenomenon; replacing spaces by their simplex categories gives categorical models with the same homotopy information.

Consider the trivial circle bundle S2 \ensuremath{\times} S1 \ensuremath{\to}S2 and the Hopf bundle (Hatcher, 2002, Examples 4.45 and 4.51)

S3 = \{(z0, z1) \ensuremath{\in}C2 : |z0|2 + |z1|2 = 1\} \ensuremath{-}\ensuremath{\to}CP1 \ensuremath{\sim}= S2, (z0, z1) 7\ensuremath{-}\ensuremath{\to}[z0 : z1]. On the chart z0 = 0, with coordinate w = z1/z0, a trivialization is

\ensuremath{\lambda}

\ensuremath{\lambda} \ensuremath{\in}S1. (w, \ensuremath{\lambda}) 7\ensuremath{-}\ensuremath{\to} 1 + |w|2 (1, w), p The other chart is analogous. Thus both maps are locally trivial bundles with fibre S1, and hence fibrations. Since the base is simply connected, their parallel transport determines the same functor from the fundamental groupoid to the homotopy category: a constant local system with value S1. Nevertheless they are not fibre homotopy equivalent. Such an equivalence would induce a weak equivalence of total spaces by the homotopy exact sequences, whereas

H1(S2 \ensuremath{\times} S1; Z) = Z.

H1(S3; Z) = 0, The missing information is the coherence of transport around two-dimensional homotopies, not an additional choice of monodromy around loops.

In contemporary language the appropriate datum is a functor from the entire homotopy type of the base, considered as an \ensuremath{\infty}-groupoid, to spaces. Passing to the fundamental groupoid of the base and the ordinary homotopy category of spaces discards part of that datum. Straightening and unstraightening provide the relevant modern comparison; Lurie (2009, Theorems 2.2.1.2 and 3.2.0.1) treat respectively space-valued and category-valued versions. Here this explains the direction of Grothendieck’s remark about \ensuremath{\infty}-groupoid stacks, without asserting that a single ordinary category HOTW (S) already determines every higher structure he envisages. The earlier package’s discussion of what a diagram theory remembers remains applicable.

There is also a small logical distinction in the definition on page 6. If the subscript 0 means that all transition functors lie in W\ensuremath{\infty}, their images lie in W once W\ensuremath{\infty}\ensuremath{\subseteq}W. The converse need not hold for a larger W. A W-locally constant diagram and an ordinarily homotopically locally constant diagram must therefore be distinguished whenever W is still variable.

\end{dcomment}

\begin{dcomment}{examples}{p8}{The walking arrow detects the failed equivalence}
Grothendieck’s first comparison theorem would identify all the categories in the large diagram after the specified localizations. The later marginal note withdraws it. The obstruction already occurs over the smallest base with a noninvertible arrow and explains why the distinction between fine and coarse equivalence occupies the rest of the chapter.

Let S = [1] = (0 \ensuremath{\to}1) and take W = W\ensuremath{\infty}. Consider the contravariant diagrams

F(1) = \ensuremath{\varnothing},

F(0) = \{\ensuremath{*}\}, G(0) = G(1) = \{\ensuremath{*}\}. There is a unique natural transformation F \ensuremath{\to}G. Its Grothendieck construction is the functor \{0\} ,\ensuremath{\to}[1] over S. Both source and target are split fibrations, and both total categories have contractible nerve, so this map is a coarse weak equivalence. Its fibre at 1 is \ensuremath{\varnothing}\ensuremath{\to}\{\ensuremath{*}\}, which is not a weak equivalence.

This difference persists after localization; it is not merely a difference between two presenting classes of arrows. Evaluation at 1 descends from the category of strict diagrams localized at objectwise weak equivalences to Hot. The displayed map has image \ensuremath{\varnothing}\ensuremath{\to}\{\ensuremath{*}\} there and therefore cannot become an isomorphism. Its coarse image is invertible by definition. Thus the forgetful comparison from fibrewise localization to coarse localization is not an equivalence. The same argument applies whenever W \ensuremath{\subseteq}W0, since then empty and nonempty categories cannot become isomorphic after applying \ensuremath{\pi}0.

The locally constant subcategory is also genuinely smaller. If a diagram H(1) \ensuremath{\to}H(0) has a weak equivalence as its transition and were isomorphic to F after objectwise localization, evaluation would make H(1) empty and H(0) contractible in Hot, contradicting the invertibility of its transition. Reversing all arrows gives the corresponding example for split opfibrations. This does not invalidate the formal localization lemma on the next pages: the missing hypothesis is that its comparison transformations lie in the particular weak-equivalence classes being used.

\end{dcomment}

\begin{dcomment}{exposition}{p12}{The localization lemma keeps its full force}
Grothendieck separates a sound formal mechanism from the geometric assertions to which he hopes to apply it. The lemma is useful precisely because its hypotheses record both kinds of information: preservation of weak equivalences by the functors, and membership of the comparison maps in the specified weak-equivalence classes.

Let P : (C, W) \ensuremath{\to}(D, V ) and R : (D, V ) \ensuremath{\to}(C, W) be relative functors. Suppose natural transformations, in either direction, connect RP with idC and PR with idD, with all components respectively in W and V . Their images under localization are natural isomorphisms. The induced functors P\textasciimacron{} and R\textasciimacron{} are therefore quasi-inverse equivalences. One needs neither a calculus of fractions nor triangle identities for the original transformations. Quasi-inverse equivalences, unlike a prescribed adjunction, require only the two natural isomorphisms.

Now suppose P factors through a full subcategory i : D0 ,\ensuremath{\to}D, written P = iP0, and put V0 = V \ensuremath{\cap}Mor(D0) and R0 = Ri. The first comparison restricts to R0P0 = RP. The second restricts to a transformation between iP0R0 and i. Because both of its endpoints lie in D0, fullness gives a unique transformation between P0R0 and idD0; faithfulness verifies its naturality. Every component belongs to V0. Consequently P0 and R0 also give inverse equivalences after localization. Applying the same argument to i shows that D0[V \ensuremath{-}1 ] \ensuremath{\to}D[V \ensuremath{-}1] is an equivalence.

This explains the apparently strong corollary about an arbitrary full subcategory containing the image of a replacement functor. The weak-equivalence class on that full subcategory must be the induced one. For example, a full subcategory of Cat/S containing all perfect replacements may calculate the coarse localization. Nothing in the lemma permits changing its induced coarse class to the fibrewise class without proving their agreement. Nor does fullness apply to the passage from splitting-preserving maps to arbitrary maps over S: that functor is generally faithful but not full. These are the two points at which the first comparison diagram asks for additional mathematics.

\end{dcomment}

\begin{dcomment}{exposition}{p15}{The precise local-to-global argument}
Grothendieck needs fibrewise equivalence to imply total equivalence without requiring a map of fibrations to preserve cartesian arrows. Smoothness or properness provides the comparison between strict fibres and comma categories that makes this possible. The earlier package’s treatment of these notions explains the lifting categories behind those comparisons; here their role is descent for a map of total categories.

For a proper functor p : X \ensuremath{\to}S, the canonical functor Xs \ensuremath{\to}X/s is a W\ensuremath{\infty}-equivalence; for a smooth functor, the corresponding functor is Xs \ensuremath{\to}s\textbackslash{}X. These assertions, in the appropriate aspherical form, are the criterion developed in Maltsiniotis (2005, Theorem 1.2 and its dual). Suppose X, Y are proper over S and f : X \ensuremath{\to}Y is any functor over S. Naturality gives

Xs

Ys

X/s

Y/s. Under W\ensuremath{\infty}\ensuremath{\subseteq}W, the vertical maps belong to W. The two-out-of-three property implies that fs \ensuremath{\in}W if and only if f/s \ensuremath{\in}W. If this holds for every s, relative L4 bis gives f \ensuremath{\in}W. For smooth functors the proof uses the undercategories and the dual relative axiom. No cartesian or cocartesian condition on f occurs.

The argument also explains the eventual universal statement on pages 41--43. Proper and smooth functors are stable under change of base. For S\ensuremath{\prime} \ensuremath{\to}S, the strict fibre of the pulled-back map at s\ensuremath{\prime} is fs\ensuremath{\prime} = fp(s\ensuremath{\prime}). Applying the same proof over S\ensuremath{\prime} shows that a fibrewise equivalence between two proper objects, or between two smooth objects, remains a total equivalence after every base change.

The parenthetical assertion here that L3 quater and L4 suffice suppresses one hypothesis: the version of Theorem A actually used is relative L4 bis. Absolute L4 directly covers the case Y = S, since S/s has a terminal object. For general Y it is the stronger formulation that is needed in the displayed proof. Grothendieck explicitly identifies this distinction on manuscript page 44. We follow that correction rather than assigning an unproved consequence to the initial axiom.

\end{dcomment}

\begin{dcomment}{exposition}{p17}{The arrow category has two different good endpoints}
Grothendieck introduces a path construction which remembers the direction of an arrow. Its endpoint maps are excellent replacement tools separately, although the pair of endpoints is not a homotopy-fibration model of the diagonal. The elementary adjunctions make all of these assertions transparent. Write Ar(A) = Fun([1], A), and let i(a) = ida, \ensuremath{\sigma}(a \ensuremath{\to}b) = a, \ensuremath{\beta}(a \ensuremath{\to}b) = b. A commutative square from i(x) to a \ensuremath{\to}b is determined by its arrow x \ensuremath{\to}a. A square from a \ensuremath{\to}b to i(y) is determined by b \ensuremath{\to}y. Therefore

\ensuremath{\beta} \ensuremath{\dashv}i \ensuremath{\dashv}\ensuremath{\sigma}.

In particular \ensuremath{\sigma}i = \ensuremath{\beta}i = id, and the squares associated with an arrow a \ensuremath{\to}b give natural transformations i\ensuremath{\sigma} \ensuremath{\Rightarrow}id and id \ensuremath{\Rightarrow}i\ensuremath{\beta}. Both endpoint projections and i are categorical homotopy equivalences, hence belong to every W satisfying L1 and L2.

The fibration properties can be checked without choosing anything. Pulling a \ensuremath{\to}b back along a\ensuremath{\prime} \ensuremath{\to}a means precomposing to obtain a\ensuremath{\prime} \ensuremath{\to}b; the resulting square is cartesian for \ensuremath{\sigma}. Pushing it forward along b \ensuremath{\to}b\ensuremath{\prime} means postcomposing to obtain a \ensuremath{\to}b\ensuremath{\prime}; that square is cocartesian for \ensuremath{\beta}. The fibres are a\textbackslash{}A and A/b, with initial and terminal identities respectively. This example already belongs to the original descent formalism; see Grothendieck (1971, Exposé VI, \S{}11(a)).

For the combined map (\ensuremath{\sigma}, \ensuremath{\beta}) : Ar(A) \ensuremath{\to}A \ensuremath{\times} A, however, the fibre at (a, b) is the discrete set HomA(a, b). Take A = [1]. Its fibre at (1, 0) is empty, while that at (0, 0) is a point, although the base is connected. Thus it is not a Hot-fibration. It is neither smooth nor proper: at (1, 0) its undercategory has the object arising from (1, 0) \ensuremath{\to}(1, 1), and its overcategory has the object arising from (0, 0) \ensuremath{\to}(1, 0), whereas its strict fibre is empty. In fact each of these comma categories is a point. The necessary fibre-to-comma weak-equivalence condition fails in both directions. This distinction is why the directed arrow construction gives \ensuremath{\Phi} and \ensuremath{\Psi}, while the two-sided replacement C requires the longer path category Ch\ensuremath{\infty}(S).

\end{dcomment}

\begin{dcomment}{exposition}{p19}{A canonical opfibration and its two projections}
Grothendieck’s replacement \ensuremath{\Psi} packages all the overcategories X/s into one category. Writing down its arrows shows both why it is strictly functorial and why a projection used to prove homotopy invariance need not be a morphism over the original base.

For p : X \ensuremath{\to}S, an object of \ensuremath{\Psi}(X) is (x, u : p(x) \ensuremath{\to}s). A morphism to (y, u\ensuremath{\prime} : p(y) \ensuremath{\to}t) is a pair

vu = u\ensuremath{\prime}p(a).

(a : x \ensuremath{\to}y, v : s \ensuremath{\to}t), The structural projection q sends this object to s and this morphism to v. For v : s \ensuremath{\to}t, the morphism (idx, v) : (x, u) \ensuremath{\to}(x, vu) is the specified cocartesian lift. Composition of these lifts is literal composition, so the opfibration is split and its fibre at s is X/s.

There is also the forgetful projection \ensuremath{\pi} : \ensuremath{\Psi}(X) \ensuremath{\to}X, which sends (x, u) to x, and the functor i : X \ensuremath{\to}\ensuremath{\Psi}(X) sending x to (x, idp(x)). A map from i(x) to (y, u\ensuremath{\prime}) is uniquely determined by a : x \ensuremath{\to}y, because its base arrow must be u\ensuremath{\prime}p(a). Hence i \ensuremath{\dashv}\ensuremath{\pi}, with \ensuremath{\pi}i = idX and counit

\ensuremath{-}\ensuremath{-}\ensuremath{-}\ensuremath{-}\ensuremath{\to}(x, u).

i\ensuremath{\pi}(x, u) = (x, idp(x)) This proves that i and \ensuremath{\pi} are inverse up to categorical homotopy. The identity qi = p makes i a functor over S, but p\ensuremath{\pi}(x, u) = p(x) generally differs from q(x, u) = s. Thus \ensuremath{\pi} generally is not a functor over S. The printed statement that i is a morphism in Cat/X is nonetheless true when \ensuremath{\Psi}(X) carries \ensuremath{\pi}: i is a section. Both descriptions of i are legitimate, and neither licenses using \ensuremath{\pi} as an inverse in Cat/S.

For a functor f : X \ensuremath{\to}Y over S, define \ensuremath{\Psi}(f)(x, u) = (f(x), u) and \ensuremath{\Psi}(f)(a, v) = (f(a), v). This preserves the specified cocartesian lifts and is strictly functorial in f. The maps i form a natural transformation in Cat/S. Their being in W shows, by the resulting commutative square, that \ensuremath{\Psi} preserves and reflects coarse weak equivalences. On fibres, however, its effect is f/s. This last observation specifies the stronger class for which \ensuremath{\Psi} becomes a replacement preserving fibrewise weak equivalences.

\end{dcomment}

\begin{dcomment}{exposition}{p20}{Which comma construction repairs which variance}
Grothendieck now isolates the extra condition that makes \ensuremath{\Psi} a relative replacement for fibrewise homotopy theory. Its fibre map is exactly f/s, so the condition is that every overcategory map lie in W. When the target is S itself this says that f is W-aspherical, since S/s has a terminal object. For a general target it is a relative condition, later called fine equivalence on the left.

The interval calculation after manuscript page 8 already shows why total equivalence cannot suffice. The complete comparison of both endpoint inclusions, given alongside manuscript page 32, also separates the two comma tests. On a smooth object, the strict fibre compares to the undercategory s\textbackslash{}X; on a proper object, to the overcategory X/s. Consequently the relevant fibrewise replacement for a fibration is \ensuremath{\Phi}, and for an opfibration it is \ensuremath{\Psi}. This is the geometric reason the later fine-right equivalence belongs on the fibration side of the large diagram, while fine-left equivalence belongs on the opfibration side. The direction names refer to the comma constructions, not to the position of the objects in that diagram.

\end{dcomment}

\begin{dcomment}{later}{p21}{From comma fibres to homotopy fibres}
Grothendieck introduces the longer path category because he now wants a replacement whose transport is homotopically invertible in both directions. The guiding theorem is Quillen’s Theorem B: if every transition X/s \ensuremath{\to}X/t induced by an arrow s \ensuremath{\to}t is a weak homotopy equivalence, then N(X/s) computes the homotopy fibre of NX \ensuremath{\to}NS over s. A precise restatement appears in Cisinski (2004, Theorem 2.3.2). This is the point at which locally constant comma data become actual fibre data for the homotopy type of a map.

For an object already Hot-fibred over S, the comparison Xs \ensuremath{\to}X/s is a weak equivalence, and the transitions between these comma categories are weak equivalences. Consequently \ensuremath{\Psi}(X) is a split opfibration with homotopically invertible transition functors. The footnote’s printed target Fibsc0S is therefore a slip: the target is Cofibsc0S. Its dual statement concerns \ensuremath{\Phi} and split fibrations. The slip does not affect the argument in the adjacent square.

The functor C uses results about Ch\ensuremath{\infty}(S) established outside the supplied pages. The input here is a natural factorization X \ensuremath{-}\ensuremath{-}\ensuremath{\to}C(X) \ensuremath{\to}S with jX \ensuremath{\in}W\ensuremath{\infty}and C(X) \ensuremath{\to}S perfect. Given this input, all subsequent uses of its functoriality are formal. For instance, the naturality square of j implies that f \ensuremath{\in}W if and only if C(f) \ensuremath{\in}W, provided W\ensuremath{\infty}\ensuremath{\subseteq}W. What is not formal is replacing the last total condition by a condition on fibres for an arbitrary larger W. That is the question on pages 22--24, and the stronger L5 introduced later addresses exactly it. The correct comparison therefore depends on which localizer is being studied, not only on the quality of the replacement object.

\end{dcomment}

\begin{dcomment}{exposition}{p23}{The contractible-base argument without a component reduction}
Grothendieck’s strategy is to inspect a chosen fibre after making the base contractible. Its force is easiest to see when the required base-change assumption is stated before the reduction. It then becomes clear that the componentwise lemma is unnecessary for this particular argument, as his marginal annotation already observes.

Assume W\ensuremath{\infty}\ensuremath{\subseteq}W, two-out-of-three, and the following additional property: pullback along a perfect functor preserves W-equivalences. Let f : X \ensuremath{\to}Y be a functor over S, with X, Y Hot-fibred over S, and suppose its total functor lies in W. Fix s \ensuremath{\in}S. Use the perfect replacement of the point \{s\} \ensuremath{\to}S to obtain

\ensuremath{-}\ensuremath{\to}jS\ensuremath{\prime} \ensuremath{-}\ensuremath{\to}rS, \{s\} j \ensuremath{\in}W\ensuremath{\infty}, where r is perfect and S\ensuremath{\prime} is weakly contractible. The distinguished object s\ensuremath{\prime} = j(s) lies over s. This construction is allowed even when S is disconnected; its image lies in the relevant component.

Put X\ensuremath{\prime} = X \ensuremath{\times}S S\ensuremath{\prime} and Y \ensuremath{\prime} = Y \ensuremath{\times}S S\ensuremath{\prime}. The extra pullback hypothesis gives f \ensuremath{\prime} : X\ensuremath{\prime} \ensuremath{\to}Y \ensuremath{\prime} in W. Base change preserves the Hot-fibration property. Since S\ensuremath{\prime} is weakly contractible, the inclusions of the fibres at s\ensuremath{\prime} into X\ensuremath{\prime} and Y \ensuremath{\prime} are weak homotopy equivalences. Equivalently, their homotopy-fibre squares are pulled back along the weak equivalence \{s\ensuremath{\prime}\} \ensuremath{\to}S\ensuremath{\prime}. In the commutative square

Xs

Ys

X\ensuremath{\prime}

Y \ensuremath{\prime}, three sides lie in W, so fs \ensuremath{\in}W. This proves the fibre-detection assertion under precisely the stated additional hypothesis. It does not prove fs \ensuremath{\in}W\ensuremath{\infty}from f \ensuremath{\prime} \ensuremath{\in}W.

The proof uses neither La nor a decomposition of f into connected components. Its genuinely substantial input is stability of W under these perfect base changes, later called L6. The speculation here that truncated homotopy equivalences should satisfy that stability is withdrawn in the later annotations; the obstruction is discussed at manuscript pages 43--45. The proof itself is valid as a conditional argument, and its dependence on L6 is the mathematical information to retain.

\end{dcomment}

\begin{dcomment}{exposition}{p24}{What the path replacement actually proves}
Grothendieck wants a category over S whose strict fibres already have the homotopy types of the homotopy fibres. The replacement jX : X \ensuremath{-}\ensuremath{\to}C(X) achieves this geometrically: jX is a weak homotopy equivalence, is natural over S, and C(X) \ensuremath{-}\ensuremath{\to}S is perfect. The difficulty is to decide which weak equivalences this replacement preserves. Two assertions must be separated, because one is formal and the other detects a substantive property of the coefficient theory.

Write W tot for maps over S whose underlying functor lies in W. Assuming W\ensuremath{\infty}\ensuremath{\subseteq}W and two-outof-three, naturality gives

f \ensuremath{\in}W tot

\ensuremath{\Leftarrow}\ensuremath{\Rightarrow}C(f) \ensuremath{\in}W tot C(f)jX = jY f, S . Thus no base-change axiom is needed to make C a functor between the coarse relative categories. If F \ensuremath{\subseteq} Cat/S is full and contains every perfect object, C factors through F. The two natural transformations j on Cat/S and its restriction to F become invertible after coarse localization. Consequently the inclusion of F induces an equivalence. This proves the assertion of manuscript page 25 under the stated elementary hypotheses, independently of the conjectural pullback assertion on the present page.

The stronger assertion is that C(f) is a W-equivalence on every fibre whenever f \ensuremath{\in}W tot S . This does require the implication later named L5 on manuscript page 43. The preceding argument with a contractible covering base explains a sufficient condition precisely. For a chosen s \ensuremath{\in}S, take a perfect S\ensuremath{\prime} \ensuremath{-}\ensuremath{\to}S with contractible total category and an object over s; it suffices to cover the component containing s. If pullback along this map preserves W, the resulting total map X\ensuremath{\prime} \ensuremath{-}\ensuremath{\to}Y \ensuremath{\prime} belongs to W. For homotopy fibrations the inclusions of the chosen fibres into X\ensuremath{\prime} and Y \ensuremath{\prime} are weak homotopy equivalences. Two-out-of-three then gives the desired fibre map in W. The final comparison gives a Wequivalence, not necessarily a weak homotopy equivalence: the sentence immediately before manuscript page 24 compresses these two conclusions.

For W = W\ensuremath{\infty}this is the familiar invariance of a homotopy fibre under equivalence of total spaces over a fixed base. For a coarser W it is an additional demand. The covering calculation at manuscript page 43 below will show its failure even when both maps to the base are discrete fibrations and discrete opfibrations. In particular the phrase “il nous faut absolument” concerns compatibility with the fibrewise localizer, not the already established compatibility with the coarse one.

\end{dcomment}

\begin{dcomment}{context}{p26}{Strict splittings and Cartesian maps}
The comparison with diagram categories depends on what the morphisms preserve. The definition on manuscript page 4 uses functors preserving the chosen Cartesian lifts strictly. A Cartesian functor between split fibrations need only carry a chosen lift to some Cartesian lift; the latter is isomorphic to the chosen lift at the image object. These conditions give strict and pseudonatural transformations, respectively. They are different conditions on actual functors, before localization.

Here is a concrete failure of fullness. Put S = [1] and let G = C2 = \{1, c\}. The projection S \ensuremath{\times}BG \ensuremath{-}\ensuremath{\to}S has its constant splitting, with chosen lift (u, 1) of u : 0 \ensuremath{-}\ensuremath{\to}1. Define an endofunctor over S which fixes the two fibre groupoids and sends a cross-fibre arrow (u, g) to (u, cg). This is a functor: composition with a vertical arrow is compatible with the formula because G is abelian. Every arrow over u is Cartesian, since its group component is invertible. Our endofunctor is therefore Cartesian, but it does not preserve the splitting, since (u, 1) is sent to (u, c). Restricting the target category to Cartesian functors does not make the forgetful functor from strictly split morphisms full.

There is a separate qualification concerning objects. The functor BC4 \ensuremath{-}\ensuremath{\to}BC2 induced by reduction modulo 2 is a Grothendieck fibration: each base arrow has a lift and every arrow is Cartesian. A splitting would be a homomorphic section C2 \ensuremath{-}\ensuremath{\to}C4. Such a section cannot exist, because each lift of the nonidentity element has order 4. Hence this fibration is not isomorphic over BC2 to an object with a splitting. This is the group-extension example already explained in Grothendieck (1971, Exposé VI, \S{}9). An equivalence over the base to a split replacement is the appropriate statement; essential surjectivity in the ordinary category Cat/S would ask for an isomorphism of categories over the base.

These distinctions do not damage the intended homotopical comparison. Equivalences of fibre categories belong to W\ensuremath{\infty}, and strict replacements can be connected by such equivalences. What needs proof is precisely that replacing coherent choices by strict ones induces the correct localized comparison. Grothendieck encounters that issue directly in the failure of the proposed unit on manuscript page 30 and returns to it on pages 55--60. The underlying classical language comes from the descent formalism of SGA 1; the new task here is to control the weak equivalences during this passage. We retain the printed assertion and read it with this necessary qualification, rather than changing the category of morphisms silently.

\end{dcomment}

\begin{dcomment}{exposition}{p30}{Why the obvious unit fails to be Cartesian}
Grothendieck has found a natural weak equivalence of total categories and now asks it to be a morphism of coherent diagrams. The request is stronger. For the fibration construction \ensuremath{\Phi}(X) the fibre at s is s\textbackslash{}X, whereas the obvious map sends x \ensuremath{\in}Xs to (x, 1s). A Cartesian lift u\ensuremath{*}y \ensuremath{-}\ensuremath{\to}y in X, lying over u : s \ensuremath{-}\ensuremath{\to}t, is sent to an arrow whose source is (u\ensuremath{*}y, 1s). The chosen Cartesian lift in \ensuremath{\Phi}(X) ending at (y, 1t) instead starts at (y, u). These two objects are connected by the lift in X, but in general they are not isomorphic in s\textbackslash{}X.

The smallest example is X = S = [1] with the identity projection. Every arrow of X is Cartesian over S. In 0\textbackslash{}[1] the two objects 10 : 0 \ensuremath{-}\ensuremath{\to}0 and u : 0 \ensuremath{-}\ensuremath{\to}1 are joined by a noninvertible arrow. Thus the image of the Cartesian arrow u under iX is not Cartesian in \ensuremath{\Phi}(X). Indeed, in the Grothendieck construction of a strict diagram, an arrow over u is Cartesian exactly when its vertical comparison component is invertible. Here that comparison component is the displayed noninvertible arrow. The obstruction remains even though both comma categories are contractible.

The square printed across manuscript pages 30--31 uses X/s, X/t and covariant symbols u\ensuremath{*}. Those are the symbols of the dual construction \ensuremath{\Psi}, while the surrounding discussion concerns \ensuremath{\Phi}. The intended argument is correct with s\textbackslash{}X, t\textbackslash{}X and contravariant transport, or after dualizing the entire discussion. This variance slip should not be turned into a purported failure of the comma construction itself.

Passing to the full category Fib0S permits the map iX, because its morphisms are arbitrary functors over S. The natural weak equivalence X \ensuremath{-}\ensuremath{\to}\ensuremath{\Phi}C(X) then proves the coarse comparison in the “théorème-remballage”. The theorem does not yet compare strict coherent diagrams by fibrewise weak equivalences. This is the exact boundary between what the argument has proved and what the next attempt must establish. The later Cartesian comparison is directed towards X, rather than away from it; its construction supplies the missing information about coherence.

\end{dcomment}

\begin{dcomment}{examples}{p32}{A relative type which is not locally constant}
Grothendieck’s reversal becomes transparent on the walking arrow. Let S = [1] = \{0 \ensuremath{-}\ensuremath{\to}1\}, let X be the one-object category over 0, and let Y = S with its identity projection. The inclusion

f : \{0\} \ensuremath{-}\ensuremath{\to}[1]

is a map of split fibrations over S. The corresponding contravariant diagrams are

X1 = \ensuremath{\varnothing}\ensuremath{-}\ensuremath{\to}X0 = \ensuremath{*},

Y1 = \ensuremath{*}\ensuremath{-}\ensuremath{\to}Y0 = \ensuremath{*}. Both total categories are contractible, so f \ensuremath{\in}W\ensuremath{\infty}. But the fibre at 1 is \ensuremath{\varnothing}\ensuremath{-}\ensuremath{\to}\ensuremath{*}, which is not a weak homotopy equivalence. Fibre evaluation is a functor on the localization by objectwise weak equivalences, and sends isomorphisms there to isomorphisms of homotopy types. Therefore the image of f in that localization is not invertible. Since its image in the coarse localization is invertible, the comparison between these two localizations is not an equivalence.

All four tests of manuscript page 41 can be computed here. For the overcategories,

X/0 = \ensuremath{*},

Y/0 = \ensuremath{*}, X/1 = \ensuremath{*}, Y/1 = [1], so f is a local, or left-fine, weak equivalence. For the undercategories,

1\textbackslash{}X = \ensuremath{\varnothing},

0\textbackslash{}X = \ensuremath{*}, 0\textbackslash{}Y = [1], 1\textbackslash{}Y = \ensuremath{*}. Thus it fails the colocal, or right-fine, test. Pulling back along the point 1 \ensuremath{-}\ensuremath{\to}S gives the same failure, so it is not a universal weak equivalence. The example lies inside the category of split fibrations, hence also inside the category of smooth functors. Its opposite supplies the corresponding example for opfibrations and proper functors. This will also resolve the contradictory sentence on manuscript page 47.

An object equivalent to a locally constant diagram must have isomorphic fibre homotopy types along every base arrow: invertibility is invariant under isomorphism in the homotopy category of diagrams. The first diagram above does not have this property. Accordingly, the full subcategory of locally constant objects is proper even when the classifying space of the indexing category is contractible. Contractibility of the base makes local systems constant up to equivalence; it does not force every coherent diagram to be locally constant.

This gives a definite negative answer to the question on manuscript page 25 for the fibrewise localization of smooth or fibred objects. It also explains why the initial thirteen-category comparison could not have been correct with a single choice of weak equivalences. A derived category of diagrams must retain the noninvertible transition \ensuremath{\varnothing}\ensuremath{-}\ensuremath{\to}\ensuremath{*}, whereas the total homotopy type in this example forgets it.

\end{dcomment}

\begin{dcomment}{examples}{p38}{What a one-sided inverse does and does not give}
Grothendieck has a useful piece of the comparison: a natural isomorphism k\textasciimacron{} \ensuremath{\psi}\textasciimacron{} \ensuremath{\simeq}1. It proves essential surjectivity of k\textasciimacron{} and faithfulness of \ensuremath{\psi}\textasciimacron{}, and it shows that \ensuremath{\psi}\textasciimacron{} is injective on isomorphism classes of objects. The statement about surjectivity on all morphism sets requires an additional hypothesis.

To locate the issue, let A be the discrete category on two objects a0, a1, and let B = \ensuremath{*}be terminal. There is a unique functor k : A \ensuremath{-}\ensuremath{\to}B; let \ensuremath{\psi} pick a0. Then k\ensuremath{\psi} = 1B strictly. Yet

HomA(a0, a1) = \ensuremath{\varnothing}\ensuremath{-}\ensuremath{\to}HomB(\ensuremath{*}, \ensuremath{*}) = \{1\}

is not surjective. Thus a functor with a section need not be full. In this example \ensuremath{\psi} is even fully faithful, so the parenthetical suggestion that failure of faithfulness of k is equivalent to failure of full faithfulness of \ensuremath{\psi} cannot be justified solely from the retract relation. For that converse take instead A = B \ensuremath{\coprod}BM, with B = \ensuremath{*}and a nontrivial monoid M, let k be the unique functor to \ensuremath{*}and let \ensuremath{\psi} select the first summand. Then \ensuremath{\psi} is fully faithful while k is not faithful.

The precise positive statement is easy and sufficient for the progress claimed so far. Choose

\ensuremath{-}\ensuremath{\to}1. For b, b\ensuremath{\prime} \ensuremath{\in}B and h : b \ensuremath{-}\ensuremath{\to}b\ensuremath{\prime}, the morphism \ensuremath{\psi}(h) maps under k, followed by conjugation \ensuremath{\varepsilon} : k\ensuremath{\psi} with \ensuremath{\varepsilon}, to h. Hence the map

HomA(\ensuremath{\psi}b, \ensuremath{\psi}b\ensuremath{\prime}) \ensuremath{-}\ensuremath{\to}HomB(b, b\ensuremath{\prime})

is split surjective. If \ensuremath{\psi}(h) = \ensuremath{\psi}(h\ensuremath{\prime}), applying k and conjugating gives h = h\ensuremath{\prime}, proving faithfulness. If \ensuremath{\psi}b \ensuremath{\simeq}\ensuremath{\psi}b\ensuremath{\prime}, the same argument gives b \ensuremath{\simeq}b\ensuremath{\prime}. None of these statements identifies an arbitrary a \ensuremath{\in}A with \ensuremath{\psi}ka.

A natural zigzag of weak equivalences between 1 and \ensuremath{\psi}k, sought on the following pages, would supply exactly that missing identification after localization. The full comparison theorem can therefore remain correct although the present formal inference is too strong. We record an error in this intermediate inference, with its corrected scope, rather than a gap in the later comparison once that zigzag has been constructed.

\end{dcomment}

\begin{dcomment}{status}{p40}{The locally constant comparison is now a theorem}
Grothendieck is asking for morphisms that retain coherent transport while forgetting arbitrary choices of Cartesian lifts. His unsuccessful unit is not evidence that the locally constant subcategory fails to embed fully. It identifies a defect of that particular representative of the comparison. The arguments on manuscript pages 55--60 repair the comparison, and pages 65--69 establish the relevant reflection.

The higher coherence retained by local constancy was illustrated at manuscript page 7. The later theory treats the local constancy question for a general derivator. Cisinski (2009, Proposition 3.10 and Theorem 3.13) proves that locally constant objects form a full subderivator and that its inclusion has both a left and a right adjoint. These are the published numbering; the earlier arXiv version numbers the corresponding statements 3.9 and 3.12. The theorem establishes the positive status of the full embedding and the existence of the two adjoints, without identifying all the coarse localizations of this chapter for an arbitrary W.

That last restriction matters. If W is homological or a finite homotopy truncation, a coarse Wequivalence of total spaces need not be a W-equivalence on homotopy fibres. The local constancy construction in a derivator still makes sense, but the additional identification with the coarse localization of Cat/S uses the stronger hypothesis isolated below. For W\ensuremath{\infty}the categorical path model and the modern description by coherent space-valued diagrams fit together; the commentary on manuscript pages 65--69 gives the precise reflector. This question is therefore settled in its intended homotopy-theoretic setting, while its unqualified extension to every localizer is false.

\end{dcomment}

\begin{dcomment}{exposition}{p41}{Four tests on a map over the base}
Grothendieck now separates the total homotopy type from the two possible diagram directions. For p : X \ensuremath{-}\ensuremath{\to}S there are three kinds of fibre to consider. The strict fibre Xs remembers only objects and arrows mapping to s and 1s. The comma category X/s also allows arrows p(x) \ensuremath{-}\ensuremath{\to}s in the base; it varies covariantly with s. The comma category s\textbackslash{}X allows arrows s \ensuremath{-}\ensuremath{\to}p(x) and varies contravariantly. Thus for f : X \ensuremath{-}\ensuremath{\to}Y over S the tests are

total

f \ensuremath{\in}W,

fibrewise fs \ensuremath{\in}W (s \ensuremath{\in}S),

local, or left-fine f/s \ensuremath{\in}W (s \ensuremath{\in}S),

colocal, or right-fine s\textbackslash{}f \ensuremath{\in}W (s \ensuremath{\in}S). The word “local” here refers to these particular comma tests, not to a chosen Grothendieck topology. The names “left” and “right” are the manuscript’s conventions. They must not be interchanged merely because the fibrations are drawn on the left side of the thirteen-category diagram.

For the ordinary weak homotopy equivalences, the role of the comma tests is to detect the maps of coherent diagrams produced by \ensuremath{\Psi} and \ensuremath{\Phi}. A map can preserve the total homotopy type while changing one of those diagrams, as the walking-arrow example at manuscript page 32 proves. Conversely, strict fibres alone do not detect a total equivalence without some condition on the maps to S. For example, the functor from the discrete two-object category \{0, 1\} to [1], over [1], is an isomorphism on both strict fibres but changes the total space from two components to a contractible space. It also fails both comma tests: over 1, and under 0, respectively, it compares a discrete two-point category to a contractible one.

The assertion concerning arbitrary base change of a fibrewise equivalence is formal: the fibre of a pullback over s\ensuremath{\prime} \ensuremath{\in}S\ensuremath{\prime} is the original fibre over its image in S. The assertion concerning base change of a coarse equivalence along every homotopy fibration is not formal for a general W; it is exactly the additional stability later called L6 bis. Footnote 17 is valid for W\ensuremath{\infty}, the setting explicitly adopted for the proposition, and must not be read as a theorem about every localizer.

The local-isomorphism formulation has an elementary explanation in the usual categorical convention of discrete fibrations. If T \ensuremath{-}\ensuremath{\to}S is a discrete fibration, then

(X \ensuremath{\times}S T)/t \ensuremath{\sim}= X/p(t).

An object on the right determines the object of T above its source by the unique lifting property. Thus local equivalence is preserved by this pullback, and relative axiom L4 bis gives total equivalence afterward. Conversely the representable discrete fibrations S/s \ensuremath{-}\ensuremath{\to}S recover the comma tests themselves. Dualizing gives the colocal statement. This explains both the strength and the orientation of the local tests without equating them to arbitrary base-change invariance.

\end{dcomment}

\begin{dcomment}{exposition}{p42}{When strict fibres compute a comma fibre}
Grothendieck’s proposition is useful because it states exactly when a strict fibre is a valid homotopical substitute for one of the comma categories. The earlier discussion of smoothness and properness in DER\_I-a, manuscript pages 52--59, supplies the geometric input; here we use that input to compare weak equivalences of objects over the same base.

If X \ensuremath{-}\ensuremath{\to}S is smooth, the canonical inclusion Xs \ensuremath{-}\ensuremath{\to}s\textbackslash{}X is a weak homotopy equivalence. When both X and Y are smooth, the commutative square

Xs

Ys

s\textbackslash{}X

s\textbackslash{}Y has vertical arrows in W if W\ensuremath{\infty}\ensuremath{\subseteq}W. Two-out-of-three proves the equivalence of the fibrewise and right-fine tests. For proper objects the same argument uses Xs \ensuremath{-}\ensuremath{\to}X/s and proves equivalence with the left-fine test. The proof requires no comparison of total spaces at this stage. The same argument works with W-smooth or W-proper objects when the vertical maps are directly known to be in W.

To obtain universality, make any change of base T \ensuremath{-}\ensuremath{\to}S. Smoothness and properness are preserved, and the fibrewise maps remain in W. The preceding comparison gives a colocal or local W-equivalence over T, respectively. Relative axiom L4 bis, or its dual, now makes the total map over T a W-equivalence. The base change was arbitrary, so the map is universally a W-equivalence. This fills the quantifier in the proof; it is not a new appeal to L6. Here base change acts on a map already known to be a fibrewise equivalence between suitably regular objects.

For homotopy fibrations both comma comparisons are weak homotopy equivalences, universally on the base. Hence the fibrewise, local and colocal conditions coincide and imply universal total equivalence. The reverse implication from a total equivalence is the extra statement later named L5. It holds for W\ensuremath{\infty} by invariance of homotopy fibres, and the next examples show why it cannot be appended for every W.

The modern extension to asphericity structures is given by Maltsiniotis (2005, Corollary 4.16 and Theorem 5.4): Grothendieck fibrations are smooth, and smoothness is characterized by the appropriate base-change property. This is the direct continuation of the present comparison argument. In highercategorical straightening one likewise distinguishes equivalence in the diagram direction from equivalence of underlying total spaces; the latter agrees with the former for locally constant fibrations, not for every right or left fibration.

\end{dcomment}

\begin{dcomment}{examples}{p43}{The circle covering as a categorical counterexample}
Grothendieck’s covering of the circle gives a complete obstruction to L5 for homology coefficients in which its degree is invertible. It is useful to realize it inside Cat, so that no passage from topological coverings to categorical fibres is left implicit.

Fix n \ensuremath{\geq}2. Let S = BZ, the one-object category with endomorphism group Z. Let X be the action groupoid for translation of Z on Z/n. Its objects are residues a, and a morphism a \ensuremath{-}\ensuremath{\to}b is an integer m with b = a + m modulo n; composition is addition. The projection p : X \ensuremath{-}\ensuremath{\to}S sends that morphism to m. Unique lifting of each integer with specified source or target makes p both a discrete fibration and a discrete opfibration. Its strict fibre is the discrete set Z/n.

The groupoid X is connected. The automorphism group of the object 0 is nZ, so the inclusion B(nZ) \ensuremath{-}\ensuremath{\to}X is an equivalence. Under this identification, p is the homomorphism nZ ,\ensuremath{\to}Z. Thus both classifying spaces are circles, and the induced map on H1(\ensuremath{-}; k) is multiplication by n after choosing the usual generators. It is the identity on H0 and both sides have zero homology in higher degrees. If k is a nonzero commutative ring in which n is invertible, p is a k-homology equivalence. Regard it now as the map from X \ensuremath{-}\ensuremath{\to}S to S \ensuremath{-}\ensuremath{\to}1S. On strict fibres it induces

H0(\ensuremath{-}; k) : kn \ensuremath{-}\ensuremath{\to}k.

Z/n \ensuremath{-}\ensuremath{\to}\ensuremath{*}, The latter map adds coordinates and has nonzero kernel. Hence it is not a fibrewise homology equivalence, although both objects over S are perfect.

The same example disproves L6 itself, not only the fibre-detection formulation. Let E be the translation groupoid of Z acting on itself, and let E \ensuremath{-}\ensuremath{\to}BZ be the evident discrete fibration and opfibration. There is exactly one arrow between each ordered pair of objects of E, so E is contractible. The objects of X \ensuremath{\times}S E are pairs (a, z); simultaneous translation preserves a \ensuremath{-}z modulo n. Each of the resulting n components has exactly one arrow between any two of its objects and is contractible. Therefore

X \ensuremath{\times}S E \ensuremath{-}\ensuremath{\to}E

induces kn \ensuremath{-}\ensuremath{\to}k on H0 and fails to be a homology equivalence. The base-change map is perfect, as required by L6. The argument explains why replacing “homotopy fibration” by “perfect” cannot repair the homological version.

\end{dcomment}

\begin{dcomment}{development}{p43}{An integral-coefficient version of the obstruction}
Grothendieck also suspects failure for integral homology, where a nontrivial finite covering of the circle no longer gives a homology equivalence. The following explicit construction settles that case. Its ingredients are standard: a semidirect product and the mapping-torus homology sequence. We checked the discussion of homological local constancy in Cisinski (2009, Scholium 1.26) and the mapping-torus calculation in Hatcher (2002, Example 2.48); we did not find this particular categorical application to the sentence here in those sources. The marking records the application and proof supplied by us, not a priority claim for the group or its homology.

Proposition (integral homology does not detect the fibre). \ours There is a perfect functor p : X \ensuremath{-}\ensuremath{\to}S which induces an isomorphism on integral homology in every degree, while its strict fibre has nonzero first integral homology. Consequently the localizer of integral homology equivalences fails L5.

Proof. Put A = Z[1/2], let h : A \ensuremath{-}\ensuremath{\to}A be multiplication by 2, and put G = A \ensuremath{\rtimes}h Z. Take X = BG, S = BZ and let p be induced by the quotient G \ensuremath{\twoheadrightarrow}Z. Since this homomorphism is surjective, every base arrow admits lifts with either prescribed source or prescribed target; all arrows in these groupoids are invertible and hence Cartesian and cocartesian. Thus p is both a Grothendieck fibration and an opfibration, and is perfect. Its strict fibre is BA.

We compute the fibre’s homology first. The group A is the filtered union of the cyclic subgroups 2\ensuremath{-}rZ. The nerve, its free simplicial chain complex and homology all commute with this filtered union: each finite tuple of group elements occurs at a finite stage, and filtered colimits of abelian groups are exact. Since BZ has the homotopy type of a circle, this gives

(i \ensuremath{\geq}2).

H0(BA; Z) = Z, H1(BA; Z) = A, Hi(BA; Z) = 0 The map induced by h acts by 2 on H1 and by the identity on H0.

The group automorphism h induces a homeomorphism of the realization of the nerve of A. Its mapping torus Th is a bundle over the circle. Its fundamental group is A \ensuremath{\rtimes}h Z after choosing the orientation of the circle, and its higher homotopy groups vanish by the homotopy sequence, since the fibre is K(A, 1). Hence Th is a classifying space for G, compatibly with the projection to BZ. The mapping-torus homology sequence has the terms

\textperiodcentered{} \textperiodcentered{} \textperiodcentered{} \ensuremath{-}\ensuremath{\to}Hi(BA)

\ensuremath{-}\ensuremath{-}\ensuremath{-}\ensuremath{\to}Hi(BA) \ensuremath{-}\ensuremath{\to}Hi(Th) \ensuremath{-}\ensuremath{\to}Hi\ensuremath{-}1(BA) \ensuremath{-}\ensuremath{\to}\textperiodcentered{} \textperiodcentered{} \textperiodcentered{} . In degree 1, 1 \ensuremath{-}h\ensuremath{*}= \ensuremath{-}1 on A is invertible; in degree 0 it is zero on Z. Exactness therefore gives Hi(Th) = 0 for i \ensuremath{\geq}2 and an isomorphism H1(Th) \ensuremath{-}\ensuremath{\to}Z. Naturality with respect to the map of fibres BA \ensuremath{-}\ensuremath{\to}\ensuremath{*}identifies this last isomorphism with the homology map induced by the projection to the circle. Both spaces are connected, so the projection is also an isomorphism on H0. This proves that p is an integral homology equivalence. Its map of fibres to the fibre of 1S is BA \ensuremath{-}\ensuremath{\to}\ensuremath{*}, which fails on H1 = A = 0. \ensuremath{\Box}

The example distinguishes absolute homology of the total category from homology after changing the base. The monodromy multiplies the fibre’s homology by 2, and 1 \ensuremath{-}2 is invertible; its contribution disappears from the total homology without disappearing from the fibre. Thus the obstruction is present even with integral coefficients, independently of inverting the degree of a finite covering.

\end{dcomment}

\begin{dcomment}{exposition}{p44}{Componentwise locality and relative Theorem A}
Grothendieck separates two axiomatic matters here. The first concerns disjoint components; the second concerns descent over a base that need not be the target of the map being tested. Keeping them separate clarifies which part of the preceding argument each axiom repairs.

The component criterion is exactly equivalent to L7. To prove the nontrivial direction, let f : X \ensuremath{-}\ensuremath{\to}Y belong to W and decompose Y = ` j Yj into its nonempty connected components. Put Xj = f \ensuremath{-}1(Yj), a union of components of X. Then f = ` j fj, so L7(a) implies fj \ensuremath{\in}W for every j. L7(b) implies that every Xj is nonempty and connected. Hence \ensuremath{\pi}0(f) is bijective, and the fj are precisely the component maps. Conversely, if \ensuremath{\pi}0(f) is bijective and every component map belongs to W, L7(a) gives f \ensuremath{\in}W. The empty target creates no exception: a functor into the empty category has empty source.

To recover L7 from the criterion, consider a family (fi). The component map of ` i fi is just a component map of one fi; moreover its map on components is bijective exactly when each \ensuremath{\pi}0(fi) is bijective. The criterion therefore gives both directions of L7(a). If the target of a W-equivalence is connected, its component bijection makes the source connected, giving L7(b). This is a complete proof of the lemma, including the reflection direction in the coproduct assertion.

For the second matter, the original L4 applies to a map X \ensuremath{-}\ensuremath{\to}Y when the overcategories of its fibres over objects of Y are aspherical. In Proposition 1, however, one compares X/s and Y/s for an independent base S. These overcategories need not be aspherical; only the map between them is assumed to be a W-equivalence. Thus the printed sentence identifying the proposition directly with the original L4 does not establish it. L4 bis is exactly the required relative formulation. The dual condition uses s\textbackslash{}X \ensuremath{-}\ensuremath{\to}s\textbackslash{}Y .

The accepted modern definition of a fundamental localizer includes the relative formulation; see Cisinski (2004, \S{}1.1). Its minimal member is W\ensuremath{\infty}by that paper’s Theorem 2.2.11. This gives a precise later framework for the corrected argument. We make no claim here that the absolute and relative formulations are logically independent under every possible package of the earlier axioms; what the text has exhibited is the need for the relative assertion in this proof. Opposite-category invariance then transfers the relative assertion to its dual, using the category-of-simplices comparison discussed earlier in the chapter.

\end{dcomment}

\begin{dcomment}{status}{p45}{Why the new fibration axiom selects full homotopy theory}
Grothendieck’s second L5 is much stronger than the first axiom with that label. The first passed from equivalences of fibres to an equivalence of totals. The new one attempts the reverse implication for homotopy fibrations. The annotations correctly recognize that this is incompatible with the proposed range of examples, but the displayed uniqueness claim needs the standing connectedness condition W \ensuremath{\subseteq}W0. Without it the maximal class of all functors satisfies L1, L3 quater, L4 bis and L5, so those four conditions alone do not force W = W\ensuremath{\infty}.

Here is the mechanism of the uniqueness assertion, retaining W\ensuremath{\infty}\ensuremath{\subseteq}W \ensuremath{\subseteq}W0, two-out-of-three, the new L5, and the path replacement already available in the chapter. Suppose first that T \ensuremath{-}\ensuremath{\to}\ensuremath{*} belongs to W. Its classifying space is connected by W \ensuremath{\subseteq}W0. For a chosen object t form the perfect replacement P = CT (\ensuremath{*}) \ensuremath{-}\ensuremath{\to}T of the inclusion of that object. Its total category is weakly contractible. Two-out-of-three implies P \ensuremath{-}\ensuremath{\to}T \ensuremath{\in}W. Applying L5 to this map between the homotopy fibrations P \ensuremath{-}\ensuremath{\to}T and T \ensuremath{-}\ensuremath{\to}1T shows that its fibre F \ensuremath{-}\ensuremath{\to}\ensuremath{*}belongs to W. The fibre models the based loop space of BT. Repeating the argument with F and successive fibres gives categorical models for all iterated loop spaces, each W-aspherical and hence connected. Therefore

\ensuremath{\pi}n(BT, t) = \ensuremath{\pi}0(\ensuremath{\Omega}nBT) = 0

(n \ensuremath{\geq}1). Together with connectedness this says that T is weakly contractible.

Now let f : X \ensuremath{-}\ensuremath{\to}Y belong to W. Replace it by the perfect map CY (X) \ensuremath{-}\ensuremath{\to}Y . Its total map is still in W, and L5 makes every strict fibre W-aspherical. What we have just proved makes those fibres weakly contractible. Since this is a homotopy fibration, its homotopy fibres are weakly contractible, so its total map is a weak homotopy equivalence; the replacement X \ensuremath{-}\ensuremath{\to}CY (X) is also one. Thus f \ensuremath{\in}W\ensuremath{\infty}. This proves W = W\ensuremath{\infty}. The argument uses the previously established homotopy-fibre property of the path replacement, not a new assertion that arbitrary strict fibres compute homotopy fibres. The later abstract characterization in Cisinski (2004, Theorem 2.3.6) likewise includes preservation of connected components.

One can see the failure for every finite truncation directly. Let W\ensuremath{\leq}n mean maps inducing isomorphisms on \ensuremath{\pi}i for 0 \ensuremath{\leq}i \ensuremath{\leq}n, with all basepoints understood. Choose a nonzero abelian group A and a categorical model S of K(A, n + 1). The perfect replacement P = CS(\ensuremath{*}) \ensuremath{-}\ensuremath{\to}S has contractible total space, so belongs to W\ensuremath{\leq}n, while its fibre models K(A, n) and is not W\ensuremath{\leq}n-aspherical. For n = 0 that fibre is the nontrivial discrete group A. This disproves L5. Pulling P \ensuremath{-}\ensuremath{\to}S back along itself yields a total space weakly equivalent to \ensuremath{\Omega}BS, because the base P is contractible and the map is a homotopy fibration. The projection P \ensuremath{\times}S P \ensuremath{-}\ensuremath{\to}P therefore fails to belong to W\ensuremath{\leq}n, disproving L6 even for this perfect base change. The finite truncations are homotopical theories, but they do not satisfy the unrestricted fibre-detection requirement of this L5.

\end{dcomment}

\begin{dcomment}{exposition}{p47}{Reading the corrected classification}
Grothendieck is now trying to separate the homotopy theory of a diagram from the homotopy type of its total category. This separation is the reason for retaining three localizations. Write WS for equivalences of total categories, W d S for the condition that every s\textbackslash{}X \ensuremath{\to}s\textbackslash{}Y belongs to W, and W g S for the analogous condition on X/s \ensuremath{\to}Y/s. The superscripts record the two comma directions, not the placement of a category in the drawing. One sentence in the present paragraph must be corrected mathematically. The assertion that WS coincides with W d S on the smooth or fibred side, and dually with W g S on the proper or cofibred side, is false. The edition prints it without a negation. On the fibred side it is fibrewise equivalence that coincides with W d S; total-category equivalence remains a weaker condition. The walking-arrow example in the commentary at manuscript p. 32 proves the distinction: \{0\} ,\ensuremath{\to}[1], viewed over [1], is an equivalence of total homotopy types, but its fibres over 1 are respectively empty and nonempty. The same paragraph subsequently says that there are two notions on each border, and p. 49 uses the correct distinction.

On the six internal categories the equality with WS uses the new L5 of p. 43. This is different from the earlier axiom carrying that name. Under L5, a map between homotopy fibrations which is an equivalence on total categories is an equivalence on fibres. The ordinary replacement arguments compare all the coarse localizations without this additional axiom, and comma resolutions independently compare the fine theories. L5 is needed when we identify the coarse class on internal objects with their fine class, and thereby realize the coarse theory as a fully faithful locally constant part of the fine theory. This distinction prevents the corrected theorem from inheriting the false equality in its provisional summary.

\end{dcomment}

\begin{dcomment}{exposition}{p50}{The covariant theory and the extra opposite}
Grothendieck wants a notation which makes the symmetry between fibrations and opfibrations immediate. There is a genuine symmetry, but the edition alternates between two incompatible formulas. With the definition fixed on manuscript p. 3, the precise identities are

[Cofibsc(S)][(W g

HOTW (S) = [Sop, Cat][W \ensuremath{-}1 S)\ensuremath{-}1] \ensuremath{\simeq}[S, Cat][W \ensuremath{-}1 obj] = HOTW (Sop). obj], Here Wobj means objectwise membership in W. There is no opposite on the value category in the last expression. The displayed isomorphism in the middle of p. 50 has this correct form; the preceding prose and the editorial correction at p. 51 insert an additional opposite. We retain those readings as part of the source, while using the comma-category models to determine the variance of every assertion.

Indeed, taking opposite total categories sends a fibration over Sop to an opfibration over S. It sends a functor X \ensuremath{\to}Y to a functor Xop \ensuremath{\to}Y op in the same direction. Thus it defines a covariant isomorphism of the categories of structured objects; it does not reverse their morphisms. The same observation justifies the correctly stated symmetry of the perfect-fibration localization at p. 51.

The extra opposite cannot be repaired by an unspecified equivalence. Take S = e and W = W\ensuremath{\infty}. The category of homotopy types is not equivalent to its opposite. Its initial object \ensuremath{\varnothing}is strict: a morphism T \ensuremath{\to}\ensuremath{\varnothing}exists only when T \ensuremath{\simeq}\ensuremath{\varnothing}. Such an equivalence would force the initial object of the opposite category, namely the one-point homotopy type, to be strict there. This would say that every map \ensuremath{*}\ensuremath{\to}T is an isomorphism, contradicted by the inclusion of one point into a two-point set. The external opposite is therefore a substantive variance slip, even though both fine localizations themselves are correctly constructed.

The notation HOTo W (S) may still be used as a label for the covariant theory, provided it is defined by the left-fine localization. It should not be interpreted here as taking the opposite of HOTW (Sop). This convention also gives the stated direction to the functor j in the diagram which follows.

\end{dcomment}

\begin{dcomment}{context}{p53}{Why the free split fibration is the right comparison}
The difficulty with the first attempt is a difficulty about morphisms, not about the existence of enough good objects. A path completion puts every object into a particularly well-behaved class, but its unit need not preserve a chosen splitting. The comma construction \ensuremath{\Phi} has a stronger universal property: it freely supplies the selected Cartesian lifts. That property gives the reverse comparison map with exactly the morphism convention required by Fibsc(S).

This construction was already part of the theory of fibred categories before the present manuscript. Gray’s treatment constructs the free split fibration, and the modern formulation identifies split fibrations and splitting-preserving functors with strict algebras for the comma-category 2-monad on Cat/S (Gray, 1966; Emmenegger et al., 2024). In the latter reference, Theorem 2.6 separates strict algebras from pseudo-algebras; Theorem 5.4 treats the two adjoints used to split an already fibred category. These are distinct operations. The present \ensuremath{\Phi} is the free construction on an arbitrary category over S.

Its importance for localization is elementary. Let U : Fibsc(S) \ensuremath{\to}Cat/S forget the splitting. There are transformations

\ensuremath{\eta} : id \ensuremath{-}\ensuremath{\to}U\ensuremath{\Phi},

\ensuremath{\gamma} : \ensuremath{\Phi}U \ensuremath{-}\ensuremath{\to}id. The first lives in Cat/S; the second lives in the category with splitting-preserving morphisms. Both become isomorphisms after the relevant localization. Consequently U and \ensuremath{\Phi} induce inverse equivalences. One does not need a comparison transformation in a category with an incorrectly enlarged or restricted class of arrows.

For coarse equivalences, this conclusion requires only that W contain categorical homotopy equivalences and satisfy two-out-of-three. The ordinary functor forgetting the initial base object gives U\ensuremath{\Phi}(X) \ensuremath{\to}X, with a section \ensuremath{\eta}X and a natural contraction. Although the forgetful map is usually not over S, the contraction proves that \ensuremath{\eta}X is an equivalence of total homotopy types. The counit is then a coarse equivalence because its composite with the unit on UY is the identity. The more informative fibrewise verification comes on pp. 58--60.

\end{dcomment}

\begin{dcomment}{exposition}{p56}{The free split adjunction, with its arrows specified}
Grothendieck’s proposed adjunction is the efficient way to finish the comparison, once “Cartesian” in the Hom-set on the left is read as preserving the selected Cartesian lifts. We give its construction because the same calculation resolves the naturality issue in the preceding paragraph.

For p : X \ensuremath{\to}S, an object of \ensuremath{\Phi}X is (s, x, u) with u : s \ensuremath{\to}p(x). A morphism

(s, x, u) \ensuremath{-}\ensuremath{\to}(t, y, v)

is a pair (a, f) with a : s \ensuremath{\to}t, f : x \ensuremath{\to}y, and p(f)u = va. The chosen lift along a with target (t, y, v) is

\ensuremath{-}\ensuremath{-}\ensuremath{-}\ensuremath{-}\ensuremath{\to}(t, y, v).

(s, y, va) Identity and composition laws for these chosen lifts hold strictly. The unit sends x to (p(x), x, id). Let q : Y \ensuremath{\to}S be split, and let h : X \ensuremath{\to}UY be any functor over S. Write cu,z : u\ensuremath{*}z \ensuremath{\to}z for the chosen lift. Define

bh(s, x, u) = u\ensuremath{*}h(x).

For (a, f) as above, define bh(a, f) as the unique arrow over a such that

cv,h(y) bh(a, f) = h(f)cu,h(x).

The Cartesian property of cv,h(y) gives existence and uniqueness. These same uniqueness statements prove the identity and composition laws: both possible images of a composite have the required base arrow and the same composite with the terminal chosen lift. The strict splitting laws show that a selected lift (a, id) is taken to the selected lift in Y . Thus bh is splitting-preserving.

Conversely, every splitting-preserving H : \ensuremath{\Phi}X \ensuremath{\to}Y is forced by h = H\ensuremath{\eta}X. The chosen arrow (s, x, u) \ensuremath{\to}(p(x), x, id) forces its object value to be u\ensuremath{*}h(x), and Cartesian uniqueness forces its value on every morphism. The two constructions are inverse and natural, proving the actual adjunction

Homsplit/S(\ensuremath{\Phi}X, Y ) \ensuremath{\sim}= HomS(X, UY ).

Its counit is \ensuremath{\gamma}Y (s, y, u) = u\ensuremath{*}y, and \ensuremath{\gamma}Y \ensuremath{\eta}UY = idUY . Its naturality is strict for splitting-preserving functors, which are precisely the morphisms needed on this side of the comparison. This is the free-algebra adjunction underlying the 2-monad description cited above.

For a merely cloven fibration the same formula still gives a Cartesian functor \ensuremath{\Phi}X \ensuremath{\to}X, using the coherence isomorphisms of the cleavage. Different choices yield canonically isomorphic functors compatible with their Cartesian comparison arrows. But an arbitrary functor over S need not preserve Cartesian arrows, so the formula is not strictly natural on the full category Fib(S). Nor can “Cartesian” alone imply preservation of the selected splitting. Those stronger assertions in the source are unnecessary; the strict adjunction just proved gives the comparison in Fibsc(S), and the next block explains what survives after localization for all functors over S.

\end{dcomment}

\begin{dcomment}{development}{p57}{Naturality recovered after localization}
Grothendieck wants the comparison to accommodate all functors over the base, since the categories Fib(S) and Liss(S) are full subcategories of Cat/S. The raw pullback construction does not have that naturality. The localization nevertheless supplies exactly the required repair. We spell out the following elementary elaboration of his argument; the “ours” label concerns this formulation and proof, not priority. We checked the free-fibration and strictification formulations in Emmenegger et al. (2024, Sections 2 and 5). They distinguish strict and pseudo preservation of cleavages but do not state the following repair for arbitrary functors over the base. We make no broader claim about its absence from the literature.

Localized naturality of pullback retractions. \ours Suppose W\ensuremath{\infty}\ensuremath{\subseteq}W. Choose normalized cleavages on fibrations X \ensuremath{\to}S. Then the functors \ensuremath{\gamma}X : \ensuremath{\Phi}X \ensuremath{\to}X determine a natural isomorphism between the functors induced by U\ensuremath{\Phi} and the identity on the right-fine localization of the full category Fib(S). The same conclusion holds for coarse localization.

Proof. First, any natural transformation over S between functors F, G : A \ensuremath{\to}B makes their images equal in (Cat/S)[(W d S)\ensuremath{-}1]. It is represented by a functor H : A \ensuremath{\times} [1] \ensuremath{\to}B over S. The projection r : A \ensuremath{\times} [1] \ensuremath{\to}A is right-fine: for every s its comma functor is the projection

(s\textbackslash{}A) \ensuremath{\times} [1] \ensuremath{-}\ensuremath{\to}s\textbackslash{}A,

a categorical homotopy equivalence. Its two sections become the same inverse of r, hence F = Hj0 and G = Hj1 become equal. The identical proof applies to coarse localization. If A is fibred, A \ensuremath{\times} [1] is fibred, so this argument also takes place within the indicated full subcategory.

Now let f : X \ensuremath{\to}Y be any functor over S, without a Cartesian hypothesis. For (s, x, u) define a vertical arrow

\ensuremath{\theta}f,(s,x,u) : f(u\ensuremath{*}

Xx) \ensuremath{-}\ensuremath{\to}u\ensuremath{*} Y f(x) by the equation

cY

u,f(x)\ensuremath{\theta}f,(s,x,u) = f(cX u,x). Existence and uniqueness follow from the Cartesian property of the lift in Y . This defines a natural transformation f\ensuremath{\gamma}X \ensuremath{\Rightarrow}\ensuremath{\gamma}Y \ensuremath{\Phi}(f) over S. To check naturality, compose the two candidate arrows with the chosen lift to the terminal object of a morphism (a, g) in \ensuremath{\Phi}X. Both composites are f(g)f(cX u,x), u,x) = f(gcX and both have the same base arrow; Cartesian uniqueness identifies them. The first paragraph of the proof therefore gives naturality after localization.

The comparison maps are right-fine equivalences. On a fibre, \ensuremath{\gamma}X,s is right adjoint to Xs ,\ensuremath{\to}s\textbackslash{}X: an arrow y \ensuremath{\to}x over u : s \ensuremath{\to}p(x) factors uniquely through u\ensuremath{*} Xx. Every adjunction gives a homotopy equivalence of nerves. Since source and target of \ensuremath{\gamma}X are fibrations, right-fine equivalence is equivalent to fibrewise equivalence. Thus the localized natural transformation is an isomorphism. Its coarse version follows either from the same argument or from \ensuremath{\gamma}X\ensuremath{\eta}X = id and coarse invertibility of \ensuremath{\eta}X. \ensuremath{\Box}

The defect before localization is real. Take S = [1], X = S, Y = S \ensuremath{\times} [1], and f(s) = (s, s). At the object (0, 1, 0 \ensuremath{\to}1) of \ensuremath{\Phi}X, the two functors in the comparison square have values (0, 0) and (0, 1). These objects are not isomorphic. The comparison \ensuremath{\theta} is the noninvertible arrow between them. Thus the source’s strict naturality assertion cannot simply be replaced by a claim of naturality up to isomorphism; the localization, or the restriction to splitting-preserving morphisms, performs an essential part of the repair.

\end{dcomment}

\begin{dcomment}{exposition}{p59}{The comma resolution is right-fine}
Grothendieck has reduced the fine comparison to checking that the unit and counit are fine equivalences. This is the decisive point at which the comma categories preserve the relative diagram, rather than merely its total homotopy type. A direct contraction proves the assertion for every X \ensuremath{\to}S and makes the proposed base-change verification unnecessary.

Fix s \ensuremath{\in}S. An object of s\textbackslash{}\ensuremath{\Phi}X consists of

\ensuremath{-}\ensuremath{\to}vt

\ensuremath{-}\ensuremath{\to}up(x). s Define R : s\textbackslash{}\ensuremath{\Phi}X \ensuremath{\to}s\textbackslash{}X by composition, R(v, u, x) = (x, uv). A morphism represented by (a, f) satisfies av = v\ensuremath{\prime} and p(f)u = u\ensuremath{\prime}a; it is sent to f, which satisfies p(f)uv = u\ensuremath{\prime}v\ensuremath{\prime}. Thus R is a functor. The functor J = s\textbackslash{}\ensuremath{\eta}X sends (x, w) to the factorization

\ensuremath{-}\ensuremath{\to}p(x)

\ensuremath{-}\ensuremath{\to}p(x). s We have RJ = id. The arrow (u, idx) gives a natural transformation

ids\textbackslash{}\ensuremath{\Phi}X \ensuremath{-}\ensuremath{\to}JR.

It is over s, since uv is the composite arrow of the target factorization. These identities prove that J and R are inverse up to categorical homotopy. Hence \ensuremath{\eta}X belongs to W d S whenever W contains W\ensuremath{\infty}.

For a split fibration Y , the fibre of the counit is

(y, u) 7\ensuremath{-}\ensuremath{\to}u\ensuremath{*}y.

\ensuremath{\gamma}Y,s : s\textbackslash{}Y \ensuremath{-}\ensuremath{\to}Ys, If K : Ys ,\ensuremath{\to}s\textbackslash{}Y is the inclusion, Cartesian factorization gives

Homs\textbackslash{}Y (Kz, (y, u)) \ensuremath{\sim}= HomYs(z, u\ensuremath{*}y).

Thus K \ensuremath{\dashv}\ensuremath{\gamma}Y,s, with identity unit for a normalized splitting. In particular both functors are W\ensuremath{\infty}equivalences. The counit is therefore fibrewise, and hence right-fine, an equivalence. Finally, \ensuremath{\Phi}(f)s is literally s\textbackslash{}f; this verifies compatibility of \ensuremath{\Phi} with the two localizers without an additional argument. All parts of the comparison now fit together. The unit and counit induce inverse natural isomorphisms between the localized composites. We obtain

S)\ensuremath{-}1] \ensuremath{\simeq}Fibsc(S)[W \ensuremath{-}1

(Cat/S)[(W d obj]. If F is any full subcategory of Cat/S containing the fibrations, the resolution U\ensuremath{\Phi} lands in F and its unit is a morphism of F on restriction. The same argument proves the asserted equivalence with F[(W d S)\ensuremath{-}1]. Reversing the comma directions gives the corresponding statement for opfibrations and W g S.

\end{dcomment}

\begin{dcomment}{exposition}{p60}{The hypotheses of the fine comparison}
The comparison just established has survived the earlier changes of axioms. Its proof uses the inclusion W\ensuremath{\infty}\ensuremath{\subseteq}W, the elementary saturation properties, and the smooth-fibre criterion already checked on p. 42. In particular it does not use the new L5. On split fibrations, right-fine equivalences are exactly objectwise W-equivalences of the associated strict diagrams; the analogous statement holds on split opfibrations. This is the part of the construction which remains useful for homological and truncated equivalences.

The pure coarse comparison with the internal categories also needs no new L5. The path replacement C produces homotopy fibrations, its unit X \ensuremath{\to}C(X) is a total homotopy equivalence, and naturality together with two-out-of-three makes C preserve coarse equivalences. L5 enters in the additional identification of coarse equivalences between these replacements with fibrewise, hence fine, equivalences. This is what makes C induce the functor back into the fine category used for the reflection. The later annotation correctly says that this requirement is highly restrictive. The commentary on pp. 43--45 explains why it must not be silently transferred to arbitrary homological localizers.

For W = W\ensuremath{\infty}, the modern literature gives a precise model for the coarse part: small categories, nerve weak equivalences, and smooth proper functors form a category of fibrant objects; perfect objects over a base suffice to compute the corresponding localization. See Cisinski (2009, Proposition 4.4 and Corollary 4.5). This places Grothendieck’s path replacement inside a developed homotopical framework. The explicit comma proof above is still useful, since it shows exactly which comparison remains valid when the coefficient equivalences are enlarged.

\end{dcomment}

\begin{dcomment}{exposition}{p64}{What the section of localization still has to prove}
Grothendieck has now produced the same locally constant category from both comma directions. The identity of Cat/S induces

S)\ensuremath{-}1] \ensuremath{-}\ensuremath{\to}(Cat/S)[W \ensuremath{-}1

\ensuremath{\rho} : (Cat/S)[(W d S ], and path replacement induces a functor in the reverse direction. The equation \ensuremath{\rho}i \ensuremath{\simeq}id says that the coarse category is a retract. It does not yet say that i is fully faithful; this is the distinction already encountered at manuscript p. 38.

The required additional datum is now visible. There is a natural arrow X \ensuremath{\to}C(X) in the fine category, and it becomes invertible in the coarse category. Its direction suggests a unit for \ensuremath{\rho} \ensuremath{\dashv}i, not for the opposite adjunction. The inverse coarse arrow C(X) \ensuremath{\to}X then supplies a counit. Grothendieck’s change of orientation on p. 67 is consequently forced by where the unit is invertible.

A localization with a fully faithful right adjoint is a reflection onto a full subcategory. In this case its image consists of the relative homotopy types whose transport along every base arrow is invertible. The adjective “locally constant” therefore describes an intrinsic condition on a diagram, not constancy of a chosen presentation of its total category. The path construction supplies a representative satisfying that condition and a universal map to it. The argument on the next pages establishes that universality through the two triangle identities, so that no fullness conclusion is extracted merely from a retraction.

\end{dcomment}

\begin{dcomment}{exposition}{p67}{A precise reflection lemma}
The remaining argument is formal once the two classes of arrows and the two uses of the unit are distinguished. We state the lemma in the form needed here; it also shows that the iterated-unit equality can be established without further geometry. This is the ordinary categorical form of the recognition principle for reflective localizations; compare Lurie (2009, Proposition 5.2.7.4) for its higher categorical formulation.

Lemma. Let \ensuremath{\Sigma}0 \ensuremath{\subseteq}\ensuremath{\Sigma} be classes of arrows of a category A, and suppose the localizations below exist. Let C : A \ensuremath{\to}A and \ensuremath{\eta} : id \ensuremath{\to}C satisfy

C(\ensuremath{\Sigma}) \ensuremath{\subseteq}\ensuremath{\Sigma}0,

\ensuremath{\eta}X \ensuremath{\in}\ensuremath{\Sigma}, \ensuremath{\eta}CX \ensuremath{\in}\ensuremath{\Sigma}0 (X \ensuremath{\in}A). Then the identity and C induce functors

P : A0 = A[\ensuremath{\Sigma}\ensuremath{-}1

0 ] \ensuremath{-}\ensuremath{\to}A1 = A[\ensuremath{\Sigma}\ensuremath{-}1], R : A1 \ensuremath{-}\ensuremath{\to}A0, with P \ensuremath{\dashv}R and R fully faithful. The unit is induced by \ensuremath{\eta}, and the counit is the inverse of its image in A1.

Since C(\ensuremath{\Sigma}0) \ensuremath{\subseteq}\ensuremath{\Sigma}0, the endofunctor C and its natural transformation descend to an Proof. endofunctor T = RP and a transformation e : idA0 \ensuremath{\to}T. Naturality extends across the formal inverses because both functors invert the required arrows. At every object x, both

a = eT x : Tx \ensuremath{\to}T 2x,

b = T(ex) : Tx \ensuremath{\to}T 2x are isomorphisms: the first by the third hypothesis, the second by the first two. We claim that a = b.

Naturality at ex gives bex = aex. Apply T and cancel the isomorphism b on the right to obtain T(b) = T(a). Naturality at a gives

T(a)a = eT 2xa,

so T(a) = eT 2x. Naturality at b gives

T(b)a = eT 2xb.

Substitute the preceding equalities and cancel the isomorphism T(a). This proves a = b. Notice that we have never cancelled ex, which need not be invertible in A0.

Let \ensuremath{\epsilon} : PR \ensuremath{\to}idA1 be the inverse of the transformation induced by \ensuremath{\eta} in A1. Its first triangle identity is \ensuremath{\epsilon}P xP(ex) = idP x, true by construction. For the second, every object of A1 is represented by an object X of A. On its image under R, the composite R(\ensuremath{\epsilon}X)eRX is b\ensuremath{-}1a, which is the identity by the calculation above. Hence P \ensuremath{\dashv}R. The counit is an isomorphism, so R is fully faithful. \ensuremath{\Box}

The full image is exactly the class of x for which ex is invertible. Objects in the image satisfy this by the third hypothesis; conversely, invertibility identifies x with RPx. Applied to \ensuremath{\Sigma}0 = W d S and \ensuremath{\Sigma} = WS, this proves the desired full faithfulness of i. The left-comma version proves the corresponding assertion for j, with the covariant convention explained at p. 50. Grothendieck’s concatenation argument is a particularly transparent geometric verification of the same identity, and remains useful for understanding the replacement itself.

\end{dcomment}

\begin{dcomment}{exposition}{p69}{Concatenation and the two iterated units}
Grothendieck’s path model makes the reflection tangible. An object of C(X) is an object x together with a path from p(x) to its new base point. Iterating C adds a second path. The two units place the empty path in different positions, so there is no reason for the resulting functors to be literally equal before localization.

Using the finite-path version with strict empty-path units, concatenation gives a functor over S

(x, c, c\ensuremath{\prime}) 7\ensuremath{-}\ensuremath{\to}(x, c\ensuremath{\prime}c).

\ensuremath{\mu}X : C2(X) \ensuremath{-}\ensuremath{\to}C(X), On morphisms it concatenates the two compatible path diagrams; compatibility of their endpoints makes this well-defined, and composition is respected because it is respected in each path segment. The two empty-path laws give the strict equations

\ensuremath{\mu}X\ensuremath{\eta}C(X) = idC(X),

\ensuremath{\mu}XC(\ensuremath{\eta}X) = idC(X). The relevant homotopical input is C(WS) \ensuremath{\subseteq}W d S. Since \ensuremath{\eta}X \ensuremath{\in}WS, we have C(\ensuremath{\eta}X) \ensuremath{\in}W d S. Two-out-ofthree applied to the second equation yields \ensuremath{\mu}X \ensuremath{\in}W d S. After fine localization both \ensuremath{\eta}C(X) and C(\ensuremath{\eta}X) are therefore the inverse of the same arrow \ensuremath{\mu}X, which proves their equality.

This proof does not cancel the unit X \ensuremath{\to}C(X) in the fine category. That would erase precisely the distinction the chapter has worked to establish. Instead it cancels a map between already completed objects. Nor does the argument require the associativity of concatenation: the two unit identities and two-out-of-three suffice for the triangle equation. Associativity equips C with the familiar monadic structure, but the adjunction needs only the part displayed here.

The formal lemma in the preceding block shows that the equality also follows from the stated hypotheses on pp. 65--66. Thus there is no outstanding coherence gap in this part of the proof. The choice of finite paths supplies an explicit common inverse and exhibits why adjoining paths becomes idempotent after the appropriate localization.

\end{dcomment}

\begin{dcomment}{later}{p69}{Path completion and the modern category of local systems}
The theorem reaches its most familiar modern form when W = W\ensuremath{\infty}. Let S denote the infinitycategory of spaces, and let

\ensuremath{\ell}: N(S)op \ensuremath{-}\ensuremath{\to}\ensuremath{\Pi}\ensuremath{\infty}N(S) be groupoid completion; we identify the completion of N(S)op with that of N(S) by reversing paths. Then the fine theory models homotopy coherent contravariant space-valued diagrams, and the locally constant theory models diagrams which carry every base arrow to an equivalence:

HOTW\ensuremath{\infty}(S) \ensuremath{\simeq}Ho Fun(N(S)op, S),

HOTlc

W\ensuremath{\infty}(S) \ensuremath{\simeq}Ho Fun(\ensuremath{\Pi}\ensuremath{\infty}N(S), S). The first comparison can be justified through an actual model equivalence. Thomason’s adjunction

c Sd2 : sSet \ensuremath{\rightleftarrows}Cat : Ex2 N

is a Quillen equivalence for nerve weak equivalences (Thomason, 1980, Theorem 4.9 and the introduction). Apply both functors objectwise to Sop-diagrams with their projective model structures. The right adjoint preserves fibrations and trivial fibrations objectwise. If F is projectively cofibrant and G projectively fibrant, the criterion for an adjoint pair of maps to be weak equivalences is checked at each object using the original Quillen equivalence; projective cofibrancy implies objectwise cofibrancy. Thus the diagram adjunction is again a Quillen equivalence. The simplicial-set diagram model then presents the coherent functor category by Lurie (2009, Proposition 4.2.4.4). This explains why strict category-valued diagrams, after objectwise localization, retain coherent space-valued diagrams.

Here Ho is taken after forming the category of coherent diagrams. It is not the category of ordinary functors into Ho(S). The higher transport information discussed earlier in the chapter is retained by this order of operations.

Restriction \ensuremath{\ell}\ensuremath{*}is fully faithful by the universal property of groupoid completion. Its left adjoint is left Kan extension \ensuremath{\ell}!, and the reflector on diagrams is \ensuremath{\ell}\ensuremath{*}\ensuremath{\ell}!. This identifies the universal property of Grothendieck’s C, rather than asserting that the particular path-category construction is literally a chosen infinity-categorical Kan extension. The standard straightening comparison and the adjunction for Kan extension are given in Lurie (2009, Theorem 2.2.1.2 and Proposition 4.3.3.7); the reflection for a left proper combinatorial model category of coefficients is Cisinski (2009, Corollary 2.5).

A local system of spaces on the infinity-groupoid \ensuremath{\Pi}\ensuremath{\infty}N(S) can equivalently be viewed as a space over |N(S)|, by taking its total space and homotopy fibres. Consequently the coarse theory has the interpretation

HOTlc

W\ensuremath{\infty}(S) \ensuremath{\simeq}Ho(S/|N(S)|). For a direct model-categorical comparison see Cisinski (2009, Example 1.7, equation (1.7.1)). This is a homotopy category of a slice, and not generally the ordinary slice of a homotopy category. The distinction already matters for a base with nontrivial fundamental group, as the following example shows.

These interpretations explain what became of the full-faithfulness question: for spaces it is settled by the reflector, and more general coefficient theories have their own locally constant subcategories and reflection results. They do not justify identifying an arbitrary W-coarse localization with that subcategory. Such an identification requires the additional exactness properties discussed on pp. 43--45. The existence of a locally constant reflection and the adequacy of a particular coarse test are separate assertions.

\end{dcomment}

\begin{dcomment}{examples}{p70}{Two computations of the reflection}
The reflection records how the base arrows must become invertible. For the interval S = [1], a contravariant diagram is an arrow B \ensuremath{-}\ensuremath{\to}uA. Since |N[1]| is contractible, locally constant diagrams are equivalent to constant diagrams. The reflector is

\ensuremath{-}\ensuremath{\to}uA) 7\ensuremath{-}\ensuremath{\to}(A

\ensuremath{-}\ensuremath{\to}A), (B with unit given by u at 1 and the identity at 0. To verify its universal property at the coherent level, a map from B \ensuremath{\to}A to the constant diagram on K consists of maps A \ensuremath{\to}K, B \ensuremath{\to}K, and a specified homotopy relating the second to the composite through A. The space of choices of the second map together with such a homotopy is contractible over the first. Thus the mapping space is naturally equivalent to Map(A, K), exactly as required. This argument also includes empty values.

For example, the diagram \ensuremath{\varnothing}\ensuremath{\to}\ensuremath{*}becomes the constant point diagram. Its total Grothendieck construction is the category supported at 0, while that of the reflected diagram is [1]. Their total nerves are equivalent, although the unit fails to be an objectwise equivalence at 1. This is the earlier counterexample viewed as the unit of an adjunction. The inclusion of constant diagrams also has a right adjoint, sending B \ensuremath{\to}A to B: the dual mapping-space computation gives Map(K, B). The left and right adjoints are distinct constructions even for this smallest base.

For a second computation take the one-object groupoid S = BZ. Every arrow of the base is already invertible, so every coherent diagram on S is locally constant and the reflection is an equivalence. Nevertheless the coarse slice contains more morphisms than the ordinary slice of Ho(S). Consider the object z : \ensuremath{*}\ensuremath{\to}BZ given by a base point. In the infinity-categorical slice,

EndHo(S/BZ)(z) \ensuremath{\sim}= Z.

Map/BZ(z, z) \ensuremath{\simeq}\ensuremath{\Omega}BZ, Indeed, the underlying map \ensuremath{*}\ensuremath{\to}\ensuremath{*}is unique, but being a map over BZ includes a path from the chosen base point to itself. Homotopy classes of such paths are the integers. In the ordinary category Ho(S)/BZ, the same object has only its identity endomorphism, because the over-condition is imposed merely as an equality of homotopy classes.

The associated local system has fibre Z with its regular translation action. A translation-equivariant map Z \ensuremath{\to}Z is determined by the image n of 0 and is k 7\ensuremath{\to}k + n; composition adds the integers. This recovers the endomorphism calculation directly from local systems. Path completion of the point over BZ supplies the categorical analogue of its universal-cover replacement, recovering exactly the transport that the ordinary slice would forget.

\end{dcomment}
